\documentclass[10pt,a4paper]{amsart}
\usepackage[margin=19mm]{geometry}
\usepackage{amsmath,amssymb,mathtools}
\usepackage[scr=rsfs,cal=euler]{mathalfa}
\usepackage{xfrac}
\usepackage[unicode,hidelinks,bookmarks=false]{hyperref}
\newcommand{\ie}{i.e.\ }
\newcommand{\cf}{cf.\ }
\newcommand{\resp}{respectively\ }
\newcommand{\Subsection}{\subsection}
\providecommand{\nobreakdash}{\nobreak\hbox{-}}
\setlength{\emergencystretch}{2em}
\multlinegap0pt
\usepackage{amssymb}
\newtheorem{theorem}{Theorem}
\def\B{\mathbb{B}}
\def\R{\mathbb{R}}
\newcommand{\expo}[1]{\exp\left( #1\right)}
\newcommand{\fourier}[1]{\mathcal{F}\left( #1 \right)}
\newcommand{\inner}[1]{\left\langle #1 \right\rangle}
\title{A biquaternionic reformulation of Maxwell's equations via Fourier analysis}
\author{Aar\'on Guill\'en-Villalobos, Briceyda B. Delgado, H\'ector Vargas Rodr\'iguez}
\date{Source: arXiv:2605.21412v2 (CC BY 4.0)}
\begin{document}
\maketitle

\noindent\emph{Source and licence.} A biquaternionic reformulation of Maxwell's equations via Fourier analysis. Version arXiv:2605.21412v2 (CC BY 4.0), distributed by arXiv under the Creative Commons Attribution 4.0 International licence, http://creativecommons.org/licenses/by/4.0/, as recorded in the arXiv licence metadata for this exact version and shown on the abstract page https://arxiv.org/abs/2605.21412v2. Full licence notice: \texttt{licenses/arxiv-2605.21412-CC-BY-4.0.txt}.\par\smallskip

\noindent\textit{Editorial scope.} Counted unit: the article's theorem th:right-inverse (source0.tex lines 505--510). The article's complete original proof of it is delivered with it (source0.tex lines 511--530, one \texttt{proof} environment of 20 lines). No mathematics is written, repaired or re-derived: the statement and the proof are the article's own lines, in the article's own notation, and no retained line is substituted or re-flowed.  No further source-local context is needed: every cross-reference of every retained line resolves inside the two delivered spans.  Nothing was omitted from any retained span, so the package carries no omission record.  No retained line contains a citation, so the wrapper carries no bibliography and no bibliography entry.  No retained line contains a figure environment, an image-inclusion command or drawing code, so no image is delivered and the package declares no asset dependency.  Editorial formatting only, in addition: an AMS \texttt{amsart} preamble that restates the article's own declarations and macros used by the retained lines, namely the environments \texttt{theorem} on separate counters, together with the standard packages those lines need.  The retained environments print in this document as Theorem 1 in order of appearance, whereas the article prints the same items with the numbering of its own full document.  The complete native source of the article as distributed in the arXiv e-print for this exact version is bundled unchanged as \texttt{sources/201061/source0.tex}.
\begin{theorem}\label{th:right-inverse}
    Let $w\in L^2(\mathscr{C},\B)$ then the operator $T_{\mathscr{C},\pm}[-i\text{ } \cdot]$ is a right inverse of $D\pm i\partial_t$. That is,
    \begin{equation}\label{eq:inverse-parabolic}
        (D\pm i\partial_t) T_{\mathscr{C},\pm}[\mp i \,w](x,t)= w(x,t), \qquad\forall (x,t)\in\R^3\times (0,\infty).
    \end{equation}
\end{theorem}

\begin{proof}
First, by convention, we will extend $w$ by zero on $\mathbb{R}^3\setminus \overline{\Omega}$. Notice that
    \begin{align*}
        \partial_jT_{\mathscr{C},\pm}[w](x,t)&=\partial_j\int_0^t\int_{\R^3}\expo{2\pi i\inner{x,k}}\expo{\mp2\pi(t-s)k}\fourier{w}(k,s)\, dk\, ds\\
        &=\int_0^t\int_{\R^3}2\pi ik_j\expo{2\pi i\inner{x,k}}\expo{\mp 2\pi(t-s)k}\fourier{w}(k,s)\, dk\, ds.
    \end{align*}
    Then,
    \[DT_{\mathscr{C},\pm}[w](x,t)=\int_0^t\int_{\R^3}2\pi ik\expo{2\pi i\inner{x,k}}\expo{\mp 2\pi(t-s)k}\fourier{w}(k,s)\, dk\,ds.\]
    We now compute the partial derivative with respect to $t$, applying the Leibniz integral rule.
    \begin{align*}
        \partial_t T_{\mathscr{C},\pm}[w](x,t)&=\partial_t\int_0^t\int_{\R^3}\expo{2\pi i\inner{x,k}}\expo{\mp 2\pi(t-s)k}\fourier{w}(k,s)\,dk\,ds\\
        &=\int_{\R^3}\expo{2\pi i\inner{x,k}}\fourier{w}(k,t)\, dk\\   &+\int_0^t\int_{\R^3}\partial_t\left[\expo{2\pi i\inner{x,k}}\expo{\mp 2\pi(t-s)k}\right]\fourier{w}(k,s)\, dk\,ds\\
        &=
        w(x,t)\mp \int_0^t\int_{\R^3}2\pi k\expo{2\pi i\inner{x,k}}\expo{\mp 2\pi(t-s)k}\fourier{w}(k,s)\,dk\, ds\\
        &=w(x,t)\pm iD T_{\mathscr{C},\pm}[w](x,t)(x,t).
    \end{align*}
    Therefore, we can easily verify that 
    $$(D\pm i\partial_t)T_{\mathscr{C},\pm}[w](x,t)=\pm iw(x,t), \qquad \forall (x,t)\in \R^3\times (0,\infty),$$
which in turn implies \eqref{eq:inverse-parabolic} as we desired.
\end{proof}

\end{document}
